\documentclass[12pt]{article}

\usepackage[margin=1in]{geometry}
\usepackage{setspace}
\doublespacing
\usepackage{amsmath,amssymb}
\usepackage{booktabs}
\usepackage{threeparttable}
\usepackage{caption}
\usepackage{array}
\usepackage{natbib}
\bibpunct{(}{)}{;}{a}{,}{,}
\usepackage{hyperref}
\hypersetup{colorlinks=true, linkcolor=black, citecolor=black, urlcolor=black}
\usepackage{microtype}

\title{\textbf{The Market Pricing of Proprietary AI Model Capability}}
\author{\begin{tabular}{c}
Zhuo Chen\textsuperscript{a}\quad Ruishuang Lu\textsuperscript{a}\\[0.5em]
{\footnotesize \textsuperscript{a} Center of Economic Research, Shandong University}
\end{tabular}}
\date{\small This version: June 2026}

\begin{document}
\maketitle
\thispagestyle{empty}

\begin{abstract}
\noindent
Which AI model capability measures do financial markets price? We study 60 major
AI model releases from 2024 to early 2026 and a firm-event panel of 86
AI-related listed firms. Across a specification curve, the broad Artificial
Analysis Intelligence Index has the strongest long-window signal; narrower
coding and math indexes do not. In closed-source releases, a one-standard-
deviation increase in this index is associated with a 3.0 percentage point
higher CAR$[0,{+}20]$ ($\hat{\beta}=0.0023$, wild bootstrap $p=0.008$). The
slope is weaker for open-weight releases and survives firm fixed effects,
Fama-French abnormal returns, event-level aggregation, and creator exclusions.
Markets price broad AI capability when releases preserve proprietary rents.
\end{abstract}

\medskip
\noindent\textit{JEL codes:} G14, O30, L86\\
\noindent\textit{Keywords:} artificial intelligence, large language models, event study, appropriability, abnormal returns

\newpage

\section{Introduction}
\label{sec:intro}

AI model releases now move a large part of the technology sector's information
environment. These releases differ from ordinary product announcements. A model
release discloses a new capability frontier, and that frontier can matter for
cloud demand, software margins, developer ecosystems, and competitive position.
The financial question is which part of that capability frontier equity markets
price.

The answer should depend on appropriability. A high-capability closed-source
model can support API revenue, customer lock-in, and ecosystem rents. A high-
capability open-weight model diffuses similar technical performance more broadly.
The same benchmark score can therefore imply different cash-flow news, depending
on whether the releasing ecosystem can capture the rents from the model.

This paper tests that idea using 60 major AI model release events from January
2024 through March 2026. We match releases to 86 AI-related listed firms and
estimate firm-event abnormal returns. The main outcome is the market-model
cumulative abnormal return over $[0,{+}20]$ trading days. We compare four
Artificial Analysis (AA) measures: Intelligence Index, Coding Index, Math Index,
and Media Elo. Because capability varies at the event level, we cluster
inference by event and report CR2 small-cluster corrections and Rademacher wild
cluster bootstrap $p$-values.

The metric comparison is itself part of the result. The broad AA Intelligence
Index has the strongest and most economically meaningful long-window relation
with returns. In CAR$[0,{+}20]$ specifications, 47.2 percent of AA Intelligence
estimates are significant at 5 percent and 83.3 percent are positive. Coding and
math indexes are much weaker. Media Elo has some significant specifications, but
its median long-window coefficient is slightly negative and its strongest
evidence appears in shorter windows. Thus the evidence does not say that every
AI benchmark is priced. It says that the broad benchmark that combines text,
reasoning, code, and multimodal performance is the relevant long-window pricing
signal.

The main result is concentrated in proprietary releases. For closed-source
releases, a one-point increase in the AA Intelligence Index predicts a 0.23
percentage point higher CAR$[0,{+}20]$ ($\hat{\beta}=0.0023$, CR2 $p=0.007$,
wild $p=0.008$). The index has a standard deviation of 13.1 points, so the
estimate implies a 3.0 percentage point higher abnormal return for a one-standard-
deviation capability increase. The full-sample slope is positive but only
borderline under small-cluster inference. This pattern points to a selective
capability-pricing channel rather than a general market reward for model
performance.

Open-weight releases attenuate the capability slope. In the interaction model,
the closed-source slope is 0.0023, while the interaction between intelligence and
open-weight status is -0.0037. The interaction is statistically significant under
CR2 ($p=0.048$) and weaker under the wild bootstrap ($p=0.086$). We interpret
this result as mechanism evidence, not as a stand-alone open-source effect.
At average capability, open-weight releases do not have a statistically different
mean return from closed-source releases; the difference lies in the marginal
return to additional capability.

Several checks support the closed-source result. Adding firm fixed effects leaves
the closed-source coefficient at 0.0023 and significant under event, firm, and
two-way clustering. Replacing market-model returns with Fama-French three-factor
abnormal returns keeps the coefficient positive and significant. Aggregating to
the event level removes the repeated-firm panel structure; the closed-source
event-level coefficient remains positive ($\hat{\beta}=0.0017$, $p=0.016$).
The result is also robust to excluding Google, OpenAI, Anthropic, Alibaba, Meta,
xAI, or Microsoft one at a time. We do not find a statistically significant
relation between intelligence and pre-event CARs, which reduces concern that the
main result only reflects a pre-announcement run-up.

The paper relates to three literatures. Event studies show how public
announcements reveal value-relevant information \citep{MacKinlay1997}. Work on
innovation and market value shows that technological advances are priced when
firms can capture rents from them \citep{Teece1986,HallJaffeTrajtenberg2005}.
Research on media and financial markets shows that attention and sentiment can
move prices \citep{Tetlock2007,EngelbergParsons2011}. We bring these ideas to AI
model releases by measuring technical capability directly and by showing that the
return relation is strongest for broad capability and when the release remains
proprietary.

\section{Data and Empirical Design}
\label{sec:data}

The unit of observation is a firm-event pair. We begin with 60 major AI model
release events between January 2024 and March 2026. Events enter the sample when
they are official model or model-family releases, have a confirmed public release
date, are accessible or publicly benchmarked, and can be linked to at least one
AI-related listed firm. The clean panel has 5,160 firm-event observations. The
main regression sample has 3,780 observations, 47 releases with nonmissing AA
Intelligence Index scores, and 85 listed firms with complete controls. The
closed-source sample has 2,899 observations across 36 events. The open-weight
sample has 881 observations across 11 events.

We estimate abnormal returns using a market model fitted over a pre-event
estimation window. The main outcome, CAR$[0,{+}20]$, sums abnormal returns from
the release date through 20 trading days after the release. AI releases are
clustered in calendar time, so this window should be read as medium-window
repricing in a dense news environment. The online appendix shows that the
results are similar with Fama-French three-factor abnormal returns and reports
pre-event CAR tests.

The main capability proxy is the AA Intelligence Index, a benchmark-based
measure for text, reasoning, code, and multimodal models. We use it because it is
the broadest technical capability measure in the data, not because it is the only
available benchmark. The online appendix reports the same specification-curve
exercise for the AA Coding Index, AA Math Index, and AA Media Elo. We center the
AA Intelligence Index at its sample mean, 26.6, when estimating interaction
models. Firm controls are log assets, book-to-market, prior-quarter volatility,
and momentum. All specifications include year fixed effects. Some robustness
checks add firm fixed effects.

We estimate variants of
\begin{equation}
\label{eq:main}
\text{CAR}_{ie} = \beta \textit{Capability}_{e} +
\mathbf{X}_{ie}'\gamma + \alpha_t + \varepsilon_{ie},
\end{equation}
where $i$ indexes firms and $e$ indexes release events. We estimate this
equation for each capability measure, then focus on the AA Intelligence Index
where the long-window signal is strongest. We also interact intelligence with an
open-weight indicator. Standard errors are clustered by event. The appendix
reports CR2 and wild cluster bootstrap inference because the effective variation
comes from 47 release events.

\section{Results}
\label{sec:results}

Table~\ref{tab:metrics} compares the candidate capability measures. AA
Intelligence has the largest share of significant estimates, the largest share
of positive estimates, and the largest median coefficient. The narrower coding
and math indexes do not show the same pattern. AA Media Elo is positive in some
specifications, but its median CAR$[0,{+}20]$ coefficient is slightly negative.
The rest of the paper therefore treats AA Intelligence as the primary broad-
capability measure and treats the other metrics as comparison benchmarks.

\begin{table}[htbp]
\centering
\caption{Which capability metric carries the pricing signal?}
\label{tab:metrics}
\begin{threeparttable}
\begin{tabular}{lrrrr}
\toprule
Capability measure & Specifications & 5\% significant & Positive & Median coefficient \\
\midrule
AA Intelligence Index & 36 & 47.2\% & 83.3\% & 0.000848 \\
AA Media Elo & 33 & 27.3\% & 45.5\% & -0.000085 \\
AA Coding Index & 36 & 25.0\% & 69.4\% & 0.000334 \\
AA Math Index & 33 & 9.1\% & 51.5\% & 0.000004 \\
\bottomrule
\end{tabular}
\begin{tablenotes}[flushleft]\footnotesize
\item The table summarizes CAR$[0,{+}20]$ specifications that vary the control
set, sample split, and year fixed effects. The reported shares use the clustered
standard error from each specification. The median coefficient is over all
CAR$[0,{+}20]$ specifications for the stated metric.
\end{tablenotes}
\end{threeparttable}
\end{table}

Table~\ref{tab:main} shows the core AA Intelligence estimates. Column 1 pools all capability
events and gives a positive but marginal slope under conservative inference.
Column 2 restricts the sample to closed-source releases. The coefficient rises to
0.0023 and remains significant under CR2 and wild bootstrap inference. Column 3
shows the interaction model. The closed-source slope remains positive; the
open-weight interaction is negative.

\begin{table}[htbp]
\centering
\caption{AI model capability and post-release abnormal returns}
\label{tab:main}
\begin{threeparttable}
\begin{tabular}{lccc}
\toprule
 & Full sample & Closed-source & Interaction \\
 & (1) & (2) & (3) \\
\midrule
AA Intelligence Index & 0.00152 & 0.00232 & 0.00226 \\
 & (0.00067) & (0.00062) & (0.00060) \\
CR2 $p$-value & 0.054 & 0.007 & 0.005 \\
Wild $p$-value & 0.053 & 0.008 & 0.008 \\
\addlinespace
AA Intelligence $\times$ Open-weight &  &  & -0.00373 \\
 &  &  & (0.00120) \\
CR2 $p$-value &  &  & 0.048 \\
Wild $p$-value &  &  & 0.086 \\
\addlinespace
Year fixed effects & Yes & Yes & Yes \\
Firm controls & Yes & Yes & Yes \\
Observations & 3,780 & 2,899 & 3,780 \\
Events & 47 & 36 & 47 \\
\bottomrule
\end{tabular}
\begin{tablenotes}[flushleft]\footnotesize
\item The dependent variable is market-model CAR$[0,{+}20]$. Column 3 uses the
AA Intelligence Index centered at its sample mean. Standard errors in
parentheses are clustered by release event. CR2 reports Bell-McCaffrey small-
cluster $p$-values. Wild reports Rademacher wild cluster bootstrap $p$-values
with 4,999 draws.
\end{tablenotes}
\end{threeparttable}
\end{table}

Table~\ref{tab:robust} reports the most important robustness checks. Firm fixed
effects leave the closed-source coefficient unchanged. Two-way clustering by
event and firm also leaves it significant. Using Fama-French abnormal returns
reduces the coefficient but keeps it positive. The pre-event CAR coefficient is
not statistically different from zero. At the event level, the closed-source
mean-CAR coefficient remains positive and significant. The result also survives
a rough restriction to US-traded tickers.

\begin{table}[htbp]
\centering
\caption{Robustness of the closed-source capability slope}
\label{tab:robust}
\begin{threeparttable}
\begin{tabular}{lcccc}
\toprule
Specification & Coefficient & SE & $p$-value & Events \\
\midrule
Firm fixed effects, event cluster & 0.00226 & 0.00063 & 0.001 & 36 \\
Firm fixed effects, firm cluster & 0.00226 & 0.00070 & 0.002 & 36 \\
Firm fixed effects, two-way cluster & 0.00226 & 0.00087 & 0.014 & 36 \\
Fama-French CAR$[0,{+}20]$ & 0.00118 & 0.00051 & 0.027 & 36 \\
Pre-event CAR & 0.00047 & 0.00030 & 0.119 & 36 \\
Event-level mean CAR & 0.00172 & 0.00067 & 0.016 & 36 \\
US-traded ticker screen & 0.00271 & 0.00072 & 0.001 & 36 \\
\bottomrule
\end{tabular}
\begin{tablenotes}[flushleft]\footnotesize
\item The table reports the coefficient on the AA Intelligence Index for
closed-source releases. The pre-event CAR row uses the pre-release abnormal
return as the dependent variable. The event-level row aggregates CAR$[0,{+}20]$
to one equal-weighted mean return per release. The US-traded ticker screen uses a
rule-based ticker filter and is reported as a robustness check, not as a formal
exchange-country classification.
\end{tablenotes}
\end{threeparttable}
\end{table}

Media sentiment does not explain away the capability result. Table~\ref{tab:sent}
shows that FinBERT sentiment is negatively associated with post-release returns.
In the joint model, both sentiment and capability retain their signs, although
both are marginal under CR2. One interpretation is that strongly positive media
coverage captures already anticipated releases that partly reverse after the
announcement. We treat this as a competing-information check rather than as a
fully identified sentiment mechanism.

\begin{table}[htbp]
\centering
\caption{Capability and media sentiment}
\label{tab:sent}
\begin{threeparttable}
\begin{tabular}{lccc}
\toprule
 & Sentiment only & Sentiment only & Joint \\
 & (1) & (2) & (3) \\
\midrule
FinBERT sentiment, $\pm 5$ days & -0.081 &  & -0.061 \\
CR2 $p$-value & 0.003 &  & 0.057 \\
\addlinespace
FinBERT sentiment, $\pm 20$ days &  & -0.096 &  \\
CR2 $p$-value &  & 0.015 &  \\
\addlinespace
AA Intelligence Index &  &  & 0.00158 \\
CR2 $p$-value &  &  & 0.087 \\
\addlinespace
Year fixed effects & Yes & Yes & Yes \\
Firm controls & Yes & Yes & Yes \\
Observations & 3,780 & 3,780 & 3,780 \\
Events & 47 & 47 & 47 \\
\bottomrule
\end{tabular}
\begin{tablenotes}[flushleft]\footnotesize
\item The dependent variable is market-model CAR$[0,{+}20]$. Sentiment is the
event-level mean FinBERT score. The joint model includes both sentiment and the
AA Intelligence Index. Event-level sentiment and intelligence are not
statistically related in the sample ($p=0.317$).
\end{tablenotes}
\end{threeparttable}
\end{table}

\section{Conclusion}
\label{sec:conclusion}

AI model capability is not priced uniformly across metrics or release regimes.
The strongest long-window signal comes from the broad AA Intelligence Index, not
from narrower coding or math measures. Within that broad metric, the strongest
evidence appears for closed-source releases, where a one-standard-deviation
increase predicts a 3.0 percentage point higher 20-day abnormal return.
Open-weight releases weaken that slope. The evidence is consistent with a simple
finance mechanism. Broad technical capability becomes value-relevant when firms
can capture rents from it.

The result survives firm fixed effects, alternative abnormal-return models,
event-level aggregation, and creator exclusions. It does not appear in pre-event
CARs. These checks do not make model releases random shocks. They do show that
the closed-source capability slope is not an artifact of firm composition, a
single developer, or the market-model return benchmark. In this setting, markets
appear to distinguish proprietary technical progress from AI news more broadly.

\newpage
\bibliographystyle{aer}
\begin{thebibliography}{9}

\bibitem[Artificial Analysis(2026)]{ArtificialAnalysis2026}
Artificial Analysis. 2026.
``Artificial Analysis Model Intelligence Index and Model Benchmark Data.''
Dataset and documentation, accessed in 2026.

\bibitem[Engelberg and Parsons(2011)]{EngelbergParsons2011}
Engelberg, Joseph E., and Christopher A. Parsons. 2011.
``The Causal Impact of Media in Financial Markets.''
\textit{Journal of Finance} 66(1): 67--97.

\bibitem[Fama and French(1993)]{FamaFrench1993}
Fama, Eugene F., and Kenneth R. French. 1993.
``Common Risk Factors in the Returns on Stocks and Bonds.''
\textit{Journal of Financial Economics} 33(1): 3--56.

\bibitem[Hall, Jaffe, and Trajtenberg(2005)]{HallJaffeTrajtenberg2005}
Hall, Bronwyn H., Adam Jaffe, and Manuel Trajtenberg. 2005.
``Market Value and Patent Citations.''
\textit{RAND Journal of Economics} 36(1): 16--38.

\bibitem[MacKinlay(1997)]{MacKinlay1997}
MacKinlay, A. Craig. 1997.
``Event Studies in Economics and Finance.''
\textit{Journal of Economic Literature} 35(1): 13--39.

\bibitem[Teece(1986)]{Teece1986}
Teece, David J. 1986.
``Profiting from Technological Innovation: Implications for Integration,
Collaboration, Licensing and Public Policy.''
\textit{Research Policy} 15(6): 285--305.

\bibitem[Tetlock(2007)]{Tetlock2007}
Tetlock, Paul C. 2007.
``Giving Content to Investor Sentiment: The Role of Media in the Stock Market.''
\textit{Journal of Finance} 62(3): 1139--1168.

\end{thebibliography}

\end{document}
